% Updated 30 August 2026
\documentclass[10pt]{article}
\usepackage[margin=0.62in]{geometry}
\usepackage[T1]{fontenc}
\usepackage{lmodern,microtype}
\usepackage{amsmath,amssymb,amsthm,mathtools}
\usepackage{enumitem,booktabs,tabularx,array,multicol,xcolor}
\usepackage[bookmarks=false]{hyperref}
\hypersetup{colorlinks=true,linkcolor=blue!45!black,urlcolor=blue!45!black}
\setlength{\parindent}{0pt}
\setlength{\parskip}{3pt}
\setlist{leftmargin=*,itemsep=1pt,topsep=2pt}
\newcommand{\OPT}{\operatorname{OPT}}
\newcommand{\ALG}{\operatorname{ALG}}
\newcommand{\E}{\mathbb E}
\newcommand{\Prb}{\mathbb P}
\newcommand{\R}{\mathbb R}
\newcommand{\1}{\mathbf 1}
\newtheorem*{theorem}{Theorem}
\newtheorem*{lemma}{Lemma}
\newenvironment{keybox}{\begin{quote}\small\itshape}{\end{quote}}
\newcommand{\catalogue}[4]{%
  \item \textbf{#1.} #2\par
  {\footnotesize\emph{Core mechanism:} #3
  \emph{Expected depth: #4.}}}

\begin{document}
\begin{center}
{\LARGE\bfseries Approximation Algorithms\par}
\vspace{1mm}
{Hadi Zamani \qquad Examiner: Ji\v{r}\'i Sgall\par}
{\small 30 August 2026}
\end{center}

\section*{Assigned chapters and sections}
\begin{itemize}
\item Sections 3.1--3.3
\item Chapter 5
\item Sections 6.1, 6.2, 6.5, and 13.2
\item Sections 7.1, 7.3, and 7.4
\item Sections 8.5 and 8.6
\item Section 10.2
\item Sections 16.1 and 16.3
\end{itemize}

\subsection*{Chapter/section map: main theorems and techniques}
\begingroup
\scriptsize
\setlength{\tabcolsep}{3.5pt}
\renewcommand{\arraystretch}{1.08}
\begin{center}
\begin{tabularx}{\textwidth}{@{}>{\raggedright\arraybackslash}p{0.13\textwidth}>{\raggedright\arraybackslash}X>{\raggedright\arraybackslash}X@{}}
\toprule
\textbf{Assigned source} & \textbf{Main theorem or guarantee} &
\textbf{Characteristic techniques}\\
\midrule
\S\S3.1--3.3 & Knapsack FPTAS; identical-machine makespan PTAS;
bin packing APTAS. & Scale and round numerical data; pseudo-polynomial DP;
large/small classification; configuration enumeration; linear grouping and greedy insertion.\\
Chapter 5 & Random assignment gives a $1/2$-approximation for MAX-CUT and a
$7/8$-approximation for MAX-E3SAT; LP rounding supplies expectation and
high-probability guarantees. & Linearity of expectation; conditional expectation;
independent/scaled rounding; Chernoff bounds; alteration and repair.\\
\S\S6.1--6.2 & Goemans--Williamson gives
$\alpha_{GW}\approx0.87856$ for MAX-CUT. & SDP/Gram-matrix relaxation; unit-vector
geometry; random-hyperplane rounding; angle-by-angle comparison.\\
\S\S6.5, 13.2 & A 3-colourable graph can be coloured with
$\widetilde O(n^{0.387})$ colours; stronger independent-set rounding yields
$\widetilde O(n^{1/4})$. & SDP vector colourings; repeated semicolouring;
degree-regime balancing; Gaussian threshold rounding; delete conflicts and repeat.\\
\S\S7.1, 7.3, 7.4 & Frequency-$f$ set cover has an $f$-approximation; shortest
path is recovered exactly; generalized Steiner tree/forest has a $2$-approximation.
& Grow feasible duals until constraints become tight; charge primal cost to the dual;
merge components; reverse-delete redundant edges.\\
\S\S8.5--8.6 & FRT embeds every $n$-point metric into a dominating random tree
with expected $O(\log n)$ stretch; buy-at-bulk consequently has an
$O(\log n)$ approximation. & Random hierarchical partitions; harmonic-sum stretch
analysis; solve the unique-path tree instance; map routes back using concavity.\\
\S10.2 & Maximum independent set in planar graphs has a PTAS. & Baker layering;
delete one residue class; bounded-treewidth dynamic programming; averaging over shifts.\\
\S\S16.1, 16.3 & No $(3/2-\varepsilon)$ approximation for unrelated-machine
makespan unless $P=NP$; MAX-E3SAT is hard beyond $7/8+\delta$. & Completeness/soundness
gap reductions; encode PCP verifier randomness as constraints; transfer the gap to an
approximation threshold.\\
\bottomrule
\end{tabularx}
\end{center}
\endgroup

\section*{Theorem inventory}
\begingroup
\small

\subsection*{Completeness checklist}
\noindent\textbf{Primary results.}
\S3: Thms.~3.1, 3.5--3.8, 3.12;
Ch.~5: Thms.~5.1--5.3, 5.5, 5.7, 5.10--5.12, 5.17, 5.19,
5.21, 5.23--5.24, 5.29, 5.31--5.32 and Cors.~5.18, 5.28, 5.30,
5.35;
\S6: Thms.~6.9--6.11, 6.20, 6.23--6.24 and Cors.~6.6, 6.22;
\S7: Thms.~7.1, 7.6, 7.8;
\S8: Thms.~8.17, 8.19, 8.23 and Cor.~8.20;
\S10.2: Thms.~10.9--10.12, 10.16;
\S13.2: Thm.~13.9;
\S16: Thms.~16.1--16.2, 16.18--16.19, 16.21--16.24 and
Cor.~16.20.

\noindent\textbf{Supporting results.}
Defs.~3.2--3.4, 3.9; Lemmas~3.10--3.11;
Lemmas~5.4, 5.6, 5.14--5.16, 5.20, 5.22, 5.25--5.27,
5.33--5.34 and Facts~5.8--5.9;
Def.~6.1, Fact~6.2, Lemmas~6.3, 6.7--6.8, 6.21 and
Facts~6.4--6.5;
Lemmas~7.5, 7.7; Lemmas~8.18, 8.21--8.22;
Lemmas~10.13--10.15; Lemmas~13.6--13.8.

\medskip
\noindent\emph{Depth convention.} ``Prove fully'' means that the assigned text
develops the proof and its intermediate estimates. ``Know proof architecture''
means that the text states or imports a deep theorem, or that the main exam task
is to explain the reduction/rounding pipeline rather than reproduce the imported
proof. This classification records the assigned text's treatment, not an
examiner preference.

\subsection*{Sections 3.1--3.3: rounding data and dynamic programming}
\begin{itemize}
\catalogue{Definition 3.2}{An algorithm is pseudopolynomial when its running
time is polynomial in the input length after the numeric data are encoded in
unary.}{Distinguish a number's magnitude from its binary bit length.}
{prove fully}

\catalogue{Definition 3.3}{A PTAS is a family $\{A_\varepsilon\}$ in which,
for every fixed $\varepsilon>0$, $A_\varepsilon$ is polynomial time and has
ratio $1+\varepsilon$ for minimisation or $1-\varepsilon$ for maximisation.}
{Keep $\varepsilon$ fixed when interpreting polynomial time.}{prove fully}

\catalogue{Definition 3.4}{An FPTAS is a PTAS whose running time is also
polynomial in $1/\varepsilon$.}{Check both the ordinary input length and the
accuracy parameter.}{prove fully}

\catalogue{Theorem 3.1}{Algorithm~3.1 correctly computes the optimum knapsack
value in $O(n\min\{B,V\})$ time, where $V=\sum_i v_i$.}{Induct on the item
prefix and show that the nondominated list contains a pair dominating every
feasible subset of that prefix.}{prove fully}

\catalogue{Theorem 3.5}{Algorithm~3.2 is a knapsack FPTAS: it returns value at
least $(1-\varepsilon)\OPT$ in $O(n^3/\varepsilon)$ time.}{Let
$M=\max_i v_i$, set $\mu=\varepsilon M/n$, round
$v_i'=\lfloor v_i/\mu\rfloor$, solve exactly, and charge at most
$n\mu=\varepsilon M\leq\varepsilon\OPT$.}{prove fully}

\catalogue{Theorem 3.6}{For any fixed number $m$ of identical parallel
machines, the family $\{A_k\}$ is a PTAS for makespan; $A_k$ has ratio
$1+1/k$.}{There are at most $km$ long jobs; enumerate their machine
assignments optimally and list-schedule the short jobs, whose last-job error is
at most $\OPT/k$.}{prove fully}

\catalogue{Theorem 3.7}{When the number of identical machines is part of the
input, minimum makespan still has a PTAS; for $k=\lceil1/\varepsilon\rceil$,
the target-$T$ subroutine either rejects $T$ or returns a schedule of length
at most $(1+1/k)T$.}{Round long jobs down to multiples of $T/k^2$, dynamic
program over the at most $(k+1)^{k^2}$ machine configurations, list-schedule
short jobs, and binary-search the integral target $T$.}{prove fully}

\catalogue{Theorem 3.8}{Unless $P=NP$, bin packing has no
$\rho$-approximation for any $\rho<3/2$.}{Reduce PARTITION: a yes-instance
packs in two bins, whereas a no-instance requires at least three, so a ratio
below $3/2$ distinguishes them.}{prove fully}

\catalogue{Definition 3.9}{An APTAS for minimisation is a family
$\{A_\varepsilon\}$ and a constant $c$ such that
$A_\varepsilon(I)\leq(1+\varepsilon)\OPT(I)+c$.}{Bin packing lacks the
rescaling property that would wash out the additive constant.}{prove fully}

\catalogue{Lemma 3.10}{Any packing of all pieces larger than $\gamma$ into
$\ell$ bins extends by First Fit to a packing of all pieces in at most
\[
\max\!\left\{\ell,\frac{\operatorname{SIZE}(I)}{1-\gamma}+1\right\}
\]
bins.}{If First Fit opens a new bin, every earlier bin has load at least
$1-\gamma$; sum those loads.}{prove fully}

\catalogue{Lemma 3.11}{If $I'$ is obtained from $I$ by linear grouping with
group size $k$, then
$\OPT(I')\leq\OPT(I)\leq\OPT(I')+k$, and any packing of $I'$ lifts to one of
$I$ with at most $k$ extra bins.}{Shift each rounded group to the preceding
original group and isolate the largest group of $k$ pieces.}{prove fully}

\catalogue{Theorem 3.12}{For every $\varepsilon>0$, bin packing has a
polynomial-time algorithm using at most
$(1+\varepsilon)\OPT(I)+1$ bins; hence it has an APTAS.}{Remove small pieces,
linearly group the large ones, solve the constant-type instance exactly, lift
by Lemma~3.11, and insert small pieces using Lemma~3.10.}{prove fully}
\end{itemize}

\subsection*{Chapter 5: random sampling and randomized rounding}
\begin{itemize}
\catalogue{Theorem 5.1}{For weighted MAX-SAT with nonnegative clause weights,
setting each variable true independently with probability $1/2$ is a randomized
$1/2$-approximation.}{Use clause indicators: a nonempty clause is satisfied with
probability at least $1/2$, so linearity of expectation gives
$\E[\ALG]\geq\tfrac12\sum_jw_j\geq\tfrac12\OPT$.}{prove fully}

\catalogue{Theorem 5.2}{For every constant $\varepsilon>0$, a
$(7/8+\varepsilon)$-approximation for MAX-E3SAT would imply $P=NP$.}{This is
the sharp PCP gap theorem; cross-reference Theorem~16.22 rather than duplicating
its reduction.}{know proof architecture}

\catalogue{Theorem 5.3}{For weighted MAX-CUT with nonnegative edge weights,
putting every vertex independently on either side with probability $1/2$ is a
randomized $1/2$-approximation.}{Each edge crosses with probability $1/2$;
sum the edge indicators and compare with the total edge weight.}{prove fully}

\catalogue{Theorem 5.5}{For a MAX-SAT instance with no negated unit clause,
setting every variable true independently with probability $p>1/2$ is a
randomized $\min\{p,1-p^2\}$-approximation.}{Lemma~5.4 gives the same lower
bound for every clause; sum weighted clause indicators.}{prove fully}

\catalogue{Theorem 5.7}{Weighted MAX-SAT has a randomized
$\frac12(\sqrt5-1)$-approximation.}{Complement variables so positive unit
clauses dominate; apply the biased assignment to all other clauses, use
Lemma~5.6 for the missed negative-unit weight, and choose
$p=1-p^2=(\sqrt5-1)/2$.}{prove fully}

\catalogue{Theorem 5.10}{Independent rounding of the standard MAX-SAT LP,
setting $x_i$ true with probability $y_i^*$, is a randomized
$(1-1/e)$-approximation.}{For a length-$\ell$ clause, AM--GM and its LP
constraint give satisfaction probability at least
$1-(1-z_j^*/\ell)^\ell\geq(1-1/e)z_j^*$; sum over clauses.}{prove fully}

\catalogue{Theorem 5.11}{Choosing the better of the LP-rounded assignment and
the unbiased random assignment is a randomized $3/4$-approximation for
MAX-SAT.}{Average the two expected clause contributions; for every integer
$\ell\geq1$, half of $1-(1-1/\ell)^\ell$ plus half of $1-2^{-\ell}$ is at
least $3/4$.}{prove fully}

\catalogue{Theorem 5.12}{If
$1-4^{-x}\leq f(x)\leq4^{x-1}$ on $[0,1]$, independently setting $x_i$ true
with probability $f(y_i^*)$ is a randomized $3/4$-approximation for MAX-SAT.}
{The product bound gives
$\Pr[C_j\text{ false}]\leq4^{-z_j^*}$; concavity of
$1-4^{-z}$ gives $1-4^{-z_j^*}\geq(3/4)z_j^*$.}{prove fully}

\catalogue{Theorem 5.17}{Let $(x^*,y^*)$ solve the prize-collecting
Steiner-tree LP, let $0<\gamma<1$, choose
$\alpha\sim U[\gamma,1]$, and run the Steiner subroutine on
$\{i:y_i^*\geq\alpha\}$. Then
\[
\E\!\left[\sum_{e\in T}c_e+\sum_{i\notin V(T)}\pi_i\right]
\leq \frac{2\ln(1/\gamma)}{1-\gamma}\sum_ec_ex_e^*
 \frac1{1-\gamma}\sum_i\pi_i(1-y_i^*).
\]}
{Combine the pointwise edge bound of Lemma~5.14 with threshold integration
(Lemma~5.15) and the probability that a vertex misses the threshold
(Lemma~5.16).}{prove fully}

\catalogue{Corollary 5.18}{With $\gamma=e^{-1/2}$, the preceding algorithm is a
$(1-e^{-1/2})^{-1}\approx2.54$ approximation for prize-collecting Steiner
tree.}{Balance the two coefficients in Theorem~5.17:
$2\ln(1/\gamma)/(1-\gamma)=1/(1-\gamma)$.}{prove fully}

\catalogue{Theorem 5.19}{Algorithm~5.1 is a randomized $3$-approximation for
metric uncapacitated facility location.}{Choose the next client minimizing
$v_j^*+C_j^*$ and open $i\in N(j)$ with probability $x_{ij}^*$; disjoint
neighbourhoods charge opening cost once, and the triangle inequality gives
expected assignment cost at most $2v_\ell^*+C_\ell^*$.}{prove fully}

\catalogue{Theorem 5.21}{For one nonpreemptive machine with integer processing
times, weights, and release dates, the $\alpha$-point randomized rounding
algorithm is a randomized $2$-approximation for
$\sum_jw_jC_j$.}{Condition on $X_j=x$ and apply Lemma~5.20; averaging gives
$\E\widehat C_j\leq2C_j^*$, and Lemma~5.22 identifies the relaxation optimum.}
{prove fully}

\catalogue{Theorem 5.23}{Let $X=\sum_iX_i$ for independent Bernoulli variables
and let $L\leq\mu=\E X\leq U$. For $\delta>0$,
\[
\Prb[X\geq(1+\delta)U]
<\left(\frac{e^\delta}{(1+\delta)^{1+\delta}}\right)^U,
\]
and for $0<\delta<1$,
\[
\Prb[X\leq(1-\delta)L]
<\left(\frac{e^{-\delta}}{(1-\delta)^{1-\delta}}\right)^L.
\]}
{Apply Markov to $e^{tX}$, factor its moment generating function by
independence, use $1+x\leq e^x$, and optimize $t$.}{prove fully}

\catalogue{Theorem 5.24}{The two bounds of Theorem~5.23 remain valid when the
independent variables satisfy $X_i\in\{0,a_i\}$ with $0<a_i\leq1$.}{Use
convexity to bound
$\E e^{tX_i}=1-p_i+p_ie^{ta_i}\leq
\exp(a_ip_i(e^t-1))$, then repeat the exponential-moment proof.}{prove fully}

\catalogue{Corollary 5.28}{Under Theorem~5.24's hypotheses, if
$0\leq L\leq\E X$, then for every $\delta>0$,
$\Prb[X\leq(1-\delta)L]<e^{-L\delta^2/2}$.}{Use Lemma~5.26 for
$0<\delta<1$; at $\delta=1$ use Lemma~5.27, and for $\delta>1$ the event is
empty.}{prove fully}

\catalogue{Theorem 5.29}{For independent path rounding of the fractional
integer-multicommodity-flow LP, if its optimum congestion
$W^*\geq c\ln n$ for a constant $c$, then with high probability every edge has
load at most $W^*+\sqrt{cW^*\ln n}$.}{Apply the Chernoff upper tail to the
independent path indicators for one edge with
$\delta=\sqrt{c\ln n/W^*}\leq1$, then union-bound over at most $n^2$ edges.}
{prove fully}

\catalogue{Corollary 5.30}{If $W^*\geq1$, independent path rounding has
maximum congestion $O(\log n)\,W^*$ with high probability.}{Run the same
Chernoff calculation with the upper benchmark
$U=(c\ln n)W^*$ and $\delta=1$, then union-bound.}{prove fully}

\catalogue{Theorem 5.31}{It is NP-hard to distinguish graphs that are
$3$-colourable from graphs whose chromatic number is at least $5$.}{The text
imports the colouring gap hardness and uses it only to calibrate later
algorithms.}{know proof architecture}

\catalogue{Theorem 5.32}{Assuming the stated variant of Unique Games, for every
constant $k>3$ it is NP-hard to distinguish $3$-colourable graphs from graphs
whose chromatic number is at least $k$.}{Treat the Unique-Games reduction as
the imported gap-producing step.}{know proof architecture}

\catalogue{Corollary 5.35}{Let $G$ be $\delta$-dense and include each vertex in
$S$ independently with probability $3c\ln n/(\delta n)$. With probability at
least $1-2n^{-(c-1)}$,
$|S|\leq6c\ln n/\delta$ and every $v\notin S$ has a neighbour in $S$.}
{Combine Lemmas~5.33--5.34 and union-bound over all vertices and the sample-size
failure event.}{prove fully}

\catalogue{Lemma 5.4}{If $p>1/2$ and a MAX-SAT instance has no negated unit
clause, every clause is satisfied under independent $p$-biased assignment
with probability at least $\min\{p,1-p^2\}$.}{A unit clause is positive;
for length at least two, unsatisfaction has probability
$p^a(1-p)^b\leq p^{a+b}\leq p^2$.}{prove fully}

\catalogue{Lemma 5.6}{After complementing variables so the positive unit
clause has at least the weight of the negative one, if $v_i$ is the negative
unit-clause weight (zero when absent), then
$\OPT\leq\sum_jw_j-\sum_iv_i$.}{Every assignment loses at least the smaller
weight from each opposite unit-clause pair.}{prove fully}

\catalogue{Fact 5.8}{For nonnegative $a_1,\dots,a_k$,
$(\prod_i a_i)^{1/k}\leq k^{-1}\sum_i a_i$.}{Apply the
arithmetic--geometric mean inequality in the clause-product calculation.}
{prove fully}

\catalogue{Fact 5.9}{If $f$ is concave on $[0,1]$, $f(0)=a$, and
$f(1)=a+b$, then $f(x)\geq a+bx$ for $x\in[0,1]$.}{A concave graph lies
above the chord joining its endpoints.}{prove fully}

\catalogue{Lemma 5.14}{For a fixed threshold $\alpha>0$, the Steiner tree
returned for terminals $\{i:y_i^*\geq\alpha\}$ satisfies
$\sum_{e\in T}c_e\leq(2/\alpha)\sum_ec_ex_e^*$.}{Scale the LP edge variables
by $1/\alpha$ to obtain a feasible Steiner relaxation and apply the
factor-$2$ Steiner-tree algorithm.}{prove fully}

\catalogue{Lemma 5.15}{If $\alpha\sim U[\gamma,1]$, then
$\E[\sum_{e\in T}c_e]\leq
\frac{2\ln(1/\gamma)}{1-\gamma}\sum_ec_ex_e^*$.}{Integrate the
$2/\alpha$ bound of Lemma~5.14 over $[\gamma,1]$.}{prove fully}

\catalogue{Lemma 5.16}{Under the same threshold experiment,
$\E[\sum_{i\notin V(T)}\pi_i]\leq
\frac1{1-\gamma}\sum_i\pi_i(1-y_i^*)$.}{Bound the probability that
$y_i^*<\alpha$ separately for $y_i^*\geq\gamma$ and $y_i^*<\gamma$.}
{prove fully}

\catalogue{Lemma 5.20}{In the release-date scheduling rounding,
$\E[\widehat C_j\mid X_j=x]\leq p_j+2x$.}{Bound idle time before job $j$ by
$x-\tfrac12$ and expected earlier processing by $x+\tfrac12$ using the
time-indexed capacity constraints.}{prove fully}

\catalogue{Lemma 5.22}{The explicitly constructed integral solution
$(y^*,C^*)$ is optimal for the integer relaxation (5.3).}{Use an interchange
argument for the compact preemptive schedule; the text leaves the proof as
Exercise~5.11.}{know proof architecture}

\catalogue{Lemma 5.25 (Markov)}{For every nonnegative random variable $X$ and
$a>0$, $\Prb[X\geq a]\leq\E X/a$.}{The expectation is at least
$a\Prb[X\geq a]$.}{prove fully}

\catalogue{Lemma 5.26}{For $0\leq\delta\leq1$,
$[e^\delta/(1+\delta)^{1+\delta}]^U\leq e^{-U\delta^2/3}$; for
$0\leq\delta<1$,
$[e^{-\delta}/(1-\delta)^{1-\delta}]^L\leq e^{-L\delta^2/2}$.}
{Take logarithms and compare derivatives/Taylor bounds on the stated
intervals.}{prove fully}

\catalogue{Lemma 5.27}{If the independent variables satisfy
$X_i\in\{0,a_i\}$ with $0<a_i\leq1$, $X=\sum_iX_i$, and $L\leq\E X$, then
$\Prb[X=0]<e^{-L}$.}{Independence gives $\prod_i(1-p_i)$; use
$1-p_i\leq e^{-p_i}\leq e^{-a_ip_i}$ and sum the means.}{prove fully}

\catalogue{Lemma 5.33}{With vertex-sampling probability
$3c\ln n/(\delta n)$,
$\Prb[|S|\geq6c\ln n/\delta]\leq n^{-c/\delta}$.}{Apply the Chernoff upper
tail with $\E|S|=3c\ln n/\delta$ and $\delta=1$.}{prove fully}

\catalogue{Lemma 5.34}{If $G$ is $\delta$-dense, then for each fixed
$v\notin S$,
$\Prb[N(v)\cap S=\varnothing]\leq n^{-3c}$.}{The expected sampled-neighbour
count is at least $3c\ln n$; apply Lemma~5.27.}{prove fully}
\end{itemize}

\subsection*{Sections 6.1, 6.2, 6.5, and 13.2: semidefinite rounding}
\begin{itemize}
\catalogue{Definition 6.1}{A symmetric $X\in\R^{n\times n}$ is positive
semidefinite, written $X\succeq0$, iff $y^TXy\geq0$ for every $y\in\R^n$.}
{Use the quadratic-form definition when translating between matrices and Gram
vectors.}{prove fully}

\catalogue{Fact 6.2}{For symmetric $X$, the following are equivalent:
$X\succeq0$; every eigenvalue is nonnegative; $X$ has a nonnegative orthogonal
spectral decomposition; and $X=V^TV$ for some $V$. Hence the SDP and vector
formulations have the same feasible objective values.}{Apply the spectral
theorem and take square roots of the nonnegative eigenvalues; the columns of
$V$ are the Gram vectors.}{prove fully}

\catalogue{Lemma 6.3}{For weighted MAX-CUT, the quadratic program
\[
\max_{y_i\in\{-1,1\}}\frac12\sum_{(i,j)\in E}w_{ij}(1-y_iy_j)
\]
models the cut value exactly.}{Encode the two sides by signs: the summand is
$w_{ij}$ exactly when the endpoints have opposite signs.}{prove fully}

\catalogue{Fact 6.4}{If $r$ has independent standard-normal coordinates, then
$r/\|r\|$ is uniform on the unit sphere.}{Rotational invariance of the Gaussian
law makes every direction equiprobable.}{prove fully}

\catalogue{Fact 6.5}{The projections of a standard Gaussian vector onto unit
vectors $e_1,e_2$ are independent standard normals iff $e_1\perp e_2$.}
{The pair is jointly Gaussian with covariance $e_1\cdot e_2$; jointly Gaussian
coordinates are independent exactly at covariance zero.}{prove fully}

\catalogue{Corollary 6.6}{If $r'$ is the projection of a standard Gaussian
onto a two-dimensional plane, then $r'/\|r'\|$ is uniform on the unit circle
in that plane.}{Apply Fact~6.4 inside the plane, using Fact~6.5 for orthogonal
coordinates.}{prove fully}

\catalogue{Lemma 6.7}{Random-hyperplane rounding separates vertices $i,j$ with
probability $\pi^{-1}\arccos(v_i\cdot v_j)$.}{Project the random normal to the
plane spanned by $v_i,v_j$; the separating directions occupy two sectors of
total angle $2\arccos(v_i\cdot v_j)$.}{prove fully}

\catalogue{Lemma 6.8}{For every $x\in[-1,1]$,
\[
\frac{\arccos x}{\pi}\geq0.878\,\frac{1-x}{2}.
\]}
{Minimize the ratio
$[\arccos x/\pi]/[(1-x)/2]$ by one-variable calculus.}{prove fully}

\catalogue{Theorem 6.9}{Solving the MAX-CUT vector relaxation and rounding by
a random hyperplane is a $0.878$-approximation.}{Sum the separation
probabilities from Lemma~6.7, apply Lemma~6.8 edgewise, and use that the vector
relaxation upper-bounds $\OPT$.}{prove fully}

\catalogue{Theorem 6.10}{If MAX-CUT has a constant
$\alpha$-approximation with $\alpha>16/17$, then $P=NP$.}{The text imports
Håstad's gap reduction; retain its exact threshold.}{know proof architecture}

\catalogue{Theorem 6.11}{Assuming Unique Games, there is no constant
$\alpha$-approximation for MAX-CUT with
\[
\alpha>\alpha_{\rm GW}:=
\min_{-1\leq x\leq1}
\frac{\arccos(x)/\pi}{(1-x)/2}\geq0.878
\]
unless $P=NP$.}{The imported Unique-Games reduction matches the integrality
profile of random-hyperplane rounding.}{know proof architecture}

\catalogue{Theorem 6.20}{Every $3$-colourable $n$-vertex graph can be coloured
with at most $4\sqrt n$ colours in polynomial time.}{Repeatedly colour a
vertex and its bipartite neighbourhood with three fresh colours while its
degree is at least $\sqrt n$; greedily colour the remaining graph of maximum
degree below $\sqrt n$.}{prove fully}

\catalogue{Lemma 6.21}{For every $3$-colourable graph, the vector-colouring
program (6.8) has a feasible solution with $\lambda\leq-1/2$.}{Map the three
colour classes to the vertices of an equilateral triangle on the unit circle;
different colours have inner product $\cos(2\pi/3)=-1/2$.}{prove fully}

\catalogue{Corollary 6.22}{Every $3$-colourable graph has a feasible solution
of (6.8) with $v_i\cdot v_j=-1/2$ on every edge.}{Use the explicit vectors in
Lemma~6.21, not merely an optimum with a smaller $\lambda$.}{prove fully}

\catalogue{Theorem 6.23}{For a $3$-colourable graph of maximum degree
$\Delta$, using $t=2+\log_3\Delta$ independent random hyperplanes gives a
semicolouring with $4\Delta^{\log_3 2}$ colours with probability at least
$1/2$.}{An edge survives all $t$ cuts with probability at most $3^{-t}$;
$m\leq n\Delta/2$ bounds the expected monochromatic edges by $n/18$, and
Markov gives the success probability.}{prove fully}

\catalogue{Theorem 6.24}{Balancing high-degree neighbourhood removal with
Theorem~6.23 semicolours every $3$-colourable $n$-vertex graph using
$O(n^{\log_6 2})=O(n^{0.387})$ colours with probability at least $1/2$.}
{The two colour costs are $O(n/\sigma)$ and
$O(\sigma^{\log_3 2})$; set $\sigma=n^{\log_6 3}$ so both equal
$O(n^{\log_6 2})$.}{prove fully}

\catalogue{Lemma 13.6}{For
$S(\varepsilon)=\{i:v_i\cdot r\geq\varepsilon\}$ with standard Gaussian $r$,
\[
\Prb[i\in S(\varepsilon)]=\overline\Phi(\varepsilon),
\qquad \E|S(\varepsilon)|=n\overline\Phi(\varepsilon).
\]}{Each unit-vector projection is a standard normal; sum membership
indicators.}{prove fully}

\catalogue{Lemma 13.7}{If $G$ is $3$-colourable with maximum degree
$\Delta$, then
\[
\Prb[i\notin S'(\varepsilon)\mid i\in S(\varepsilon)]
\leq\Delta\,\overline\Phi(\sqrt3\,\varepsilon).
\]}
{For an edge, decompose $v_j=-v_i/2+\sqrt3\,u/2$ with
$u\perp v_i$; simultaneous threshold exceedance forces
$u\cdot r\geq\sqrt3\varepsilon$, then union-bound over neighbours.}
{prove fully}

\catalogue{Lemma 13.8}{For $x>0$, with standard-normal density $p$,
\[
\frac{x}{1+x^2}p(x)\leq\overline\Phi(x)\leq\frac1x p(x).
\]}
{Integrate the Gaussian tail by parts and compare with
$p'(x)=-xp(x)$.}{prove fully}

\catalogue{Theorem 13.9}{Algorithm~13.2 returns, from a $3$-colourable graph
of maximum degree $\Delta$, an independent set of expected size
\[
\Omega\!\left(n\Delta^{-1/3}(\ln\Delta)^{-1/2}\right).
\]}{Take $\varepsilon=\sqrt{(2/3)\ln\Delta}$; Lemmas~13.7--13.8 retain at
least half the threshold set in expectation, and Lemma~13.6 supplies its
size.}{prove fully}
\end{itemize}

\subsection*{Sections 7.1, 7.3, and 7.4: primal--dual algorithms}
\begin{itemize}
\catalogue{Theorem 7.1}{For weighted set cover with nonnegative set costs,
Algorithm~7.1 is an $f$-approximation, where
$f=\max_i|\{j:i\in S_j\}|$ is the maximum element frequency.}{Raise an
uncovered element's dual until a set constraint becomes tight; every selected
set is tight, and each dual variable is charged through at most $f$ selected
sets.}{prove fully}

\catalogue{Theorem 7.6}{For nonnegative edge costs, Algorithm~7.4 returns an
exact shortest $s$--$t$ path.}{Grow the dual of the current $s$-component
until a cut edge becomes tight; Lemma~7.5 makes the tight edges a tree, and
each positive dual cut is crossed once by the final path, so primal and dual
values agree.}{prove fully}

\catalogue{Theorem 7.8}{For generalized Steiner tree (Steiner forest) with
nonnegative edge costs and required pairs $(s_i,t_i)$, Algorithm~7.6 is a
$2$-approximation.}{Grow all active-component moats uniformly, add tight
edges, then reverse-delete; Lemma~7.7 bounds the primal increase by twice the
dual increase.}{prove fully}

\catalogue{Lemma 7.5}{At every point in Algorithm~7.4, the tight-edge set
$F$ is a tree containing $s$.}{Each iteration adds exactly one edge leaving
the current $s$-component, so it spans a new vertex without closing a cycle.}
{prove fully}

\catalogue{Lemma 7.7}{In every iteration of Algorithm~7.6, if $\mathcal C$ is
the family of active components and $F'$ is the reverse-deleted forest, then
\[
\sum_{C\in\mathcal C}|\delta(C)\cap F'|\leq2|\mathcal C|.
\]}
{Contract current components in the forest and apply the degree-sum bound for
a forest, using reverse deletion to remove inactive leaves that would spoil
the charge.}{prove fully}
\end{itemize}

\subsection*{Sections 8.5 and 8.6: tree metrics and buy-at-bulk}
\begin{itemize}
\catalogue{Theorem 8.17}{For an $n$-point metric $(V,d)$ with
$d_{uv}\geq1$ for $u\neq v$, a randomized polynomial-time algorithm outputs a
tree metric $(V',T)$, $V\subseteq V'$, such that for all $u,v\in V$,
\[
d_{uv}\leq T_{uv},
\qquad \E T_{uv}\leq O(\log n)d_{uv}.
\]}
{Build a hierarchical random partition from one random permutation and
radius; bound the probability that $u,v$ split at each scale, and sum the
resulting harmonic series.}{prove fully}

\catalogue{Theorem 8.19}{If $(V',T)$ comes from a hierarchical cut
decomposition and $V$ is its leaf set, then in polynomial time one obtains a
tree metric $(V,T')$ satisfying
$T_{uv}\leq T'_{uv}\leq4T_{uv}$ for all $u,v\in V$.}{Contract each internal
Steiner node into a descendant leaf; level lengths show that contraction never
shrinks a leaf distance and increases it by at most four.}{prove fully}

\catalogue{Corollary 8.20}{Under Theorem~8.17's hypotheses, the randomized
algorithm may output a tree metric directly on $V$ satisfying
$d_{uv}\leq T'_{uv}$ and
$\E T'_{uv}\leq O(\log n)d_{uv}$.}{Compose Theorems~8.17 and 8.19 and absorb
the factor four into $O(\log n)$.}{prove fully}

\catalogue{Theorem 8.23}{In the undirected buy-at-bulk problem, assume
nonnegative edge lengths, source--sink demands, and a capacity cost $f$ with
$f(0)=0$ that is nondecreasing and has economies of scale (hence is
subadditive). Algorithm~8.6 is a randomized $O(\log n)$-approximation.}
{Embed the shortest-path metric into a dominating random tree, solve the
unique-path tree instance exactly, map paths back by Lemma~8.21, and compare
with the translated optimum using Lemma~8.22 and expected stretch.}{prove fully}

\catalogue{Lemma 8.18}{Every hierarchical-cut tree dominates the original
metric; if the least common ancestor of $u,v$ is at level $i$, then
$T_{uv}\leq2^{i+2}$.}{A level-$i$ cluster has diameter below $2^{i+1}$, while
separation below $\lfloor\log_2d_{uv}\rfloor$ is impossible; telescope the
tree-edge lengths.}{prove fully}

\catalogue{Lemma 8.21}{If $c'_{xy}$ is the capacity on tree edge
$(x,y)\in T'$, mapping each tree edge to a shortest $x$--$y$ path in $G$
produces a feasible solution of cost at most
\[
\sum_{(x,y)\in T'}T'_{xy}f(c'_{xy}).
\]}{Use $d_{xy}\leq T'_{xy}$ and subadditivity of $f$ when several mapped
tree paths share an original edge.}{prove fully}

\catalogue{Lemma 8.22}{Translate an optimal solution of $G$ to $T'$ by
installing each $c_e^*$ along the unique tree path between the endpoints of
$e$. Its tree cost is at least
$\sum_{(x,y)\in T'}T'_{xy}f(c'_{xy})$.}{Deleting $(x,y)$ separates exactly
the demands counted by $c'_{xy}$; every translated feasible solution needs
at least that capacity, and monotonicity of $f$ completes the comparison.}
{prove fully}
\end{itemize}

\subsection*{Section 10.2: planar maximum independent set}
\begin{itemize}
\catalogue{Theorem 10.9 (Theorem 1.4)}{For every constant
$\varepsilon>0$, unweighted maximum clique has no
$O(n^{\varepsilon-1})$ approximation unless $P=NP$; complementation transfers
the same general-graph hardness to maximum independent set.}{The text imports
the clique gap hardness and uses graph complementation for the transfer.}
{know proof architecture}

\catalogue{Theorem 10.10}{Maximum independent set on an $n$-vertex
$k$-outerplanar graph is solvable in
$O(2^{O(k)}n^2)$ time.}{Cross-reference the constructive treewidth proof in
Theorem~10.16; do not duplicate it here.}{know proof architecture}

\catalogue{Theorem 10.11}{For every fixed $\varepsilon>0$, Baker layering
gives a $(1-\varepsilon)$-approximation for maximum independent set in planar
graphs in time $O(2^{O(1/\varepsilon)}n^2)$.}{Partition planar layers modulo
$k+1$, delete each residue in turn, solve the resulting $k$-outerplanar
components exactly, and choose the best shift; one shift loses at most
$1/(k+1)\leq\varepsilon$ of an optimum.}{prove fully}

\catalogue{Theorem 10.12}{Given a tree decomposition with node set $V'$ and
treewidth $t$, maximum independent set is computable in
$O(2^{O(t)}|V'|)$ time.}{For every bag, tabulate the best solution in its
subtree for each independent subset of the bag; combine child tables while
subtracting duplicated bag vertices.}{prove fully}

\catalogue{Theorem 10.16}{Maximum independent set on an $n$-vertex
$k$-outerplanar graph is computable in $O(2^{O(k)}n^2)$ time.}{Use
Lemmas~10.13--10.15 to construct in $O(n^2)$ a width-$O(k)$ decomposition of
linear size, then invoke Theorem~10.12.}{prove fully}

\catalogue{Lemma 10.13}{Every $k$-outerplanar graph $G$ has a
$k$-outerplanar maximum-degree-$3$ graph $G'$ such that any tree decomposition
of $G'$ yields one of $G$ with no larger treewidth.}{Split each high-degree
vertex along the embedding, then replace all split copies by the original
name in every bag.}{prove fully}

\catalogue{Lemma 10.14}{Every $k$-outerplanar graph of maximum degree $3$ has
a maximum spanning forest of maximum load at most $3k$.}{Induct on outer
layers; after removing the exterior-face edges, the remaining graph is
$(k-1)$-outerplanar, and restoring a maximal acyclic subset increases load by
at most three.}{prove fully}

\catalogue{Lemma 10.15}{If a maximum-degree-$3$ graph has a maximum spanning
forest of maximum load at most $\ell$, then its treewidth is at most
$\ell+1$.}{Begin with the forest decomposition; for every nonforest edge, add
an endpoint along its fundamental-cycle path, so each bag gains at most the
edge/vertex load.}{prove fully}
\end{itemize}

\subsection*{Sections 16.1 and 16.3: hardness of approximation}
\begin{itemize}
\catalogue{Theorem 16.1}{It is NP-complete to decide whether an unrelated
parallel-machine instance has makespan at most $2$ when each job is allowed
only on machines in $M_j$ and has processing time $p_j\in\{1,2\}$ there.
Consequently, for every $\alpha<3/2$, an $\alpha$-approximation would imply
$P=NP$.}{Reduce three-dimensional matching using element jobs and type-specific
length-$2$ dummy jobs. Yes-instances have makespan $2$; no-instances have
integral makespan at least $3$.}{prove fully}

\catalogue{Theorem 16.2}{For every $\varepsilon>0$, directed edge-disjoint
paths admits no $\Omega(m^{-1/2+\varepsilon})$ approximation unless $P=NP$,
where $m$ is the number of arcs.}{Amplify the NP-complete two-pair instance
to $k=|A|^{\lceil1/\varepsilon\rceil}$ pairs arranged through copies of the
original graph; two routed amplified pairs force both original paths, while a
yes-instance routes all $k$.}{prove fully}

\catalogue{Theorem 16.18 (PCP theorem)}{There is a positive constant $k$ such
that
\[
NP\subseteq PCP_{1,\,1/2}(O(\log n),k).
\]}
{Interpret the polynomially many random strings as constant-query local
constraints; the PCP construction itself is imported.}{know proof architecture}

\catalogue{Theorem 16.19}{For all constants $\varepsilon,\delta>0$,
\[
NP\subseteq PCP_{1-\varepsilon,\,1/2+\delta}(O(\log n),3),
\]
and the verifier uses only the $3$-bit odd/even predicates on distinct
variables.}{Use the imported predicate-restricted PCP theorem as the source of
the completeness/soundness gap.}{know proof architecture}

\catalogue{Corollary 16.20}{If the odd/even constraint-satisfaction problem
has an $\alpha$-approximation for any constant $\alpha>1/2$, then $P=NP$;
the threshold is sharp because the all-zero or all-one assignment satisfies
at least half the constraints.}{Convert each verifier random string into one
constraint and compare completeness $1-\varepsilon$ with soundness
$1/2+\delta$; choose $\varepsilon,\delta$ around the claimed $\alpha$.}
{prove fully}

\catalogue{Theorem 16.21}{If MAX-E3SAT has an $\alpha$-approximation for any
constant $\alpha>13/14$, then $P=NP$.}{Replace every parity constraint by four
E3SAT clauses: a satisfied constraint satisfies four clauses and an
unsatisfied one three. This is an $L$-reduction with $a=7,b=1$, so
$\alpha$ transfers to $7\alpha-6>1/2$.}{prove fully}

\catalogue{Theorem 16.22 (Theorem 5.2)}{If MAX-E3SAT has an
$\alpha$-approximation for any constant $\alpha>7/8$, then $P=NP$.}
{Under the same four-clause gadget, the PCP gap
$1-\varepsilon$ versus $1/2+\delta$ becomes
$1-\varepsilon/4$ versus $7/8+\delta/4$; choose constants so an
$\alpha>7/8$ algorithm distinguishes the cases. Theorem~5.2 cross-references
this proof.}{prove fully}

\catalogue{Theorem 16.23}{For every constant $\delta>0$,
\[
NP\subseteq PCP_{1,\,7/8+\delta}(O(\log n),3),
\]
where each predicate is the OR of three queried bits or literal negations.}
{The direct perfect-completeness MAX-E3SAT PCP is imported; translate its
random strings to clauses.}{know proof architecture}

\catalogue{Theorem 16.24}{If MAX-2SAT has an $\alpha$-approximation for any
constant $\alpha>21/22$, then $P=NP$.}{The twelve-clause parity gadget allows
$11$ satisfied clauses when the parity constraint is true and at most $10$
when false, sending the PCP gap to the $21/22$ threshold. The text leaves the
gadget verification and the odd-parity analogue as exercises.}
{know proof architecture}
\end{itemize}

\endgroup

\section*{Constants and guarantees at a glance}
\begingroup
\small
\begin{center}
\begin{tabularx}{\textwidth}{@{}l>{\raggedright\arraybackslash}Xl@{}}
\toprule
Result & Hypothesis/algorithm & Exact guarantee or parameter\\
\midrule
3.5 & value scaling for knapsack & $(1-\varepsilon)$; $O(n^3/\varepsilon)$\\
3.6--3.7 & identical-machine makespan & $1+1/k$; PTAS\\
3.8, 3.12 & bin packing & hardness $<3/2$; APTAS
$(1+\varepsilon)\OPT+1$\\
5.1, 5.7, 5.10--5.12 & MAX-SAT rounding & $1/2$,
$(\sqrt5-1)/2$, $1-1/e$, $3/4$, $3/4$\\
5.18, 5.19, 5.21 & PCST, facility location, scheduling &
$(1-e^{-1/2})^{-1}\approx2.54$, $3$, $2$\\
5.23--5.30 & Chernoff/path rounding & exponents $U\delta^2/3$ and
$L\delta^2/2$; congestion $W^*+\sqrt{cW^*\ln n}$\\
5.31--5.35 & colouring gaps and dense sampling & $3$ versus $5$ colours;
UGC: $3$ versus $k$; sample success $1-2n^{-(c-1)}$\\
6.9--6.11 & MAX-CUT & $\alpha_{GW}\approx0.878$; hardness $16/17$;
UGC threshold $\alpha_{GW}$\\
6.20, 6.23--6.24 & 3-colourable graphs & $4\sqrt n$;
$4\Delta^{\log_3 2}$ with probability $\geq1/2$;
$O(n^{0.387})$ semicolours\\
7.1, 7.6, 7.8 & primal--dual & frequency $f$; exact; factor $2$\\
8.17--8.23 & FRT and buy-at-bulk & dominating expected $O(\log n)$ stretch;
$O(\log n)$ approximation\\
10.11 & planar independent set & $(1-\varepsilon)$; time
$O(2^{O(1/\varepsilon)}n^2)$\\
10.9--10.16 & clique hardness and bounded treewidth &
$O(n^{\varepsilon-1})$ hardness; $O(2^{O(t)}|V'|)$ DP\\
13.9 & independent set in 3-colourable graph &
$\Omega(n\Delta^{-1/3}(\ln\Delta)^{-1/2})$ in expectation\\
16.1--16.2 & direct gap reductions & unrelated makespan hardness
$<3/2$; directed paths hardness $\Omega(m^{-1/2+\varepsilon})$\\
16.18--16.24 & PCP/hardness & $PCP_{1,1/2}(O(\log n),k)$;
$1/2$, $13/14$, $7/8$, $21/22$ thresholds\\
\bottomrule
\end{tabularx}
\end{center}
\endgroup

\section{Language and proof template}

For a minimization problem, a ratio-$\rho\geq1$ algorithm satisfies
$\ALG(I)\leq\rho\OPT(I)$.  For a maximization problem it is often cleaner to
use a guarantee $0<\alpha\leq1$ and write $\ALG(I)\geq\alpha\OPT(I)$.
Do not silently switch between the conventions $\alpha$ and $1/\alpha$.

A PTAS gives a $(1+\varepsilon)$-approximation for minimization, or a
$(1-\varepsilon)$-approximation for maximization, in polynomial time for each
fixed $\varepsilon>0$.  An FPTAS is polynomial in both the input length and
$1/\varepsilon$.  An APTAS may have an additive constant, for example
$\ALG\leq(1+\varepsilon)\OPT+1$.  A pseudo-polynomial algorithm is polynomial
in numerical values rather than their binary encoding lengths.

\begin{keybox}
Every approximation proof has four logically separate parts: (i) the output is
feasible; (ii) a lower/upper benchmark bounds $\OPT$; (iii) the algorithm is compared
with that benchmark; (iv) the running time is polynomial in the required parameters.
\end{keybox}

For a relaxation, remember the direction:
\[
  \OPT_{\rm LP}\leq\OPT_{\rm int}\quad\text{for minimization},
  \qquad
  \OPT_{\rm LP}\geq\OPT_{\rm int}\quad\text{for maximization}.
\]
The integrality gap limits what can be proved by rounding that particular relaxation.

\section{Chapter 3: rounding data and dynamic programming}

\subsection{Knapsack: pseudo-polynomial DP and an FPTAS (Section 3.1)}

Items have sizes $s_i\in\mathbb Z_{>0}$, values $v_i\in\mathbb Z_{>0}$, and
capacity $B$.  A pseudo-polynomial DP stores nondominated pairs
$(\text{size},\text{value})$, or equivalently
\[
 D[i,z]=\min\{\text{size of a subset of }[i]\text{ of total value }z\}.
\]
The recurrence is
\[
 D[i,z]=\min\bigl(D[i-1,z],\ D[i-1,z-v_i]+s_i\bigr).
\]
Its running time is polynomial in $n$ and $\sum_i v_i$, hence not necessarily
polynomial in the bit length.

For the FPTAS, let $M=\max_i v_i$ among items that fit, set
\[
  \mu=\frac{\varepsilon M}{n},
  \qquad v_i'=\left\lfloor\frac{v_i}{\mu}\right\rfloor,
\]
and solve the rounded instance exactly by the DP.  If $S$ is returned and $O$ is
optimal, then
\begin{align*}
 \sum_{i\in S}v_i
 &\geq \mu\sum_{i\in S}v_i'
 \geq \mu\sum_{i\in O}v_i'
 \geq \sum_{i\in O}v_i-|O|\mu\\
 &\geq \OPT-\varepsilon M
 \geq(1-\varepsilon)\OPT,
\end{align*}
because $M\leq\OPT$.  Moreover
$\sum_i v_i'=O(n^2/\varepsilon)$, giving polynomial time in $n$ and
$1/\varepsilon$.

\textbf{Oral trap.}  Rounding the sizes downward can violate capacity.  Rounding the
values downward preserves feasibility and only loses objective value.

\subsection{Identical-machine makespan: a PTAS (Section 3.2)}

For $P||C_{\max}$, useful lower bounds are
\[
  \OPT\geq\max_jp_j,
  \qquad \OPT\geq\frac1m\sum_jp_j.
\]
For a guessed target $T$, classify jobs as large or small relative to
$\varepsilon T$.  There are at most $m/\varepsilon$ large jobs in any schedule
of makespan $T$.  Round their processing times to a constant number of sizes,
enumerate or dynamically program over feasible machine configurations, and place the
small jobs greedily on least-loaded machines.

If the rounded large jobs fit within $T+O(\varepsilon T)$, then immediately before a
small job is placed every machine has load at most the average load; the final overflow
is at most the size of one small job, at most $\varepsilon T$.  With constants chosen
beforehand this yields makespan at most $(1+\varepsilon)\OPT$.  Binary search or
enumeration over candidate values of $T$ removes the assumption that $\OPT$ is known.

\textbf{Why only a PTAS?}  The number of configurations is polynomial for fixed
$\varepsilon$, but its exponent may depend on $1/\varepsilon$.

\subsection{Bin packing: an APTAS (Section 3.3)}

Normalize bin capacity to one.  Separate small items from large items.  Sort the large
items and apply \emph{linear grouping}: partition them into groups of equal cardinality,
set aside the first (largest) group, and round every size in each later group up to the
largest size of that group.  Only constantly many rounded sizes remain, so a dynamic
program can optimize the rounded configurations.

The rounded solution is feasible for the original large items because sizes were
rounded up.  The discarded first group costs only $O(\varepsilon\OPT)+O(1)$ extra bins.
Insert small items greedily; if a new bin is opened, all previous bins have load close
to one.  The resulting guarantee is of the form
\[
  \ALG(I)\leq(1+\varepsilon)\OPT(I)+1.
\]
The additive one is why the result is an asymptotic PTAS rather than an ordinary PTAS.

\section{Chapter 5: randomized rounding of linear programs}

\subsection{Baseline random assignments}

For MAX-CUT, independently put each vertex on either side with probability $1/2$.
Every edge is cut with probability $1/2$, so
$\E[w(\delta(S))]=\tfrac12w(E)\geq\tfrac12\OPT$.

For a clause with $k$ distinct, noncomplementary literals under a uniform random
assignment,
\[
  \Prb[\text{clause satisfied}]=1-2^{-k}.
\]
Thus uniform random assignment gives a $7/8$ guarantee for MAX-E3SAT.  For clauses of
different signs and lengths, biased coins may improve the worst clause probability.

\subsection{Conditional expectation}

If $X$ is the objective and a random bit $b$ is the next choice, then
\[
 \E[X\mid\text{past}]
 =p\E[X\mid\text{past},b=1]+(1-p)\E[X\mid\text{past},b=0].
\]
At least one branch preserves the current conditional expectation.  Fix that branch
and continue.  This derandomizes the algorithm provided the conditional expectations
can be computed efficiently.

\subsection{Independent LP rounding and concentration}

Given a fractional solution $x_i\in[0,1]$, choose item $i$ independently with
probability $x_i$ (or a scaled probability $\min\{1,\lambda x_i\}$).  Linearity gives
the expected objective.  Chernoff bounds control a sum of independent rounded
variables: if $X=\sum_iX_i$ and $\mu=\E X$, then
\[
 \Prb[X\geq(1+\delta)\mu]
 \leq\left(\frac{e^\delta}{(1+\delta)^{1+\delta}}\right)^\mu
 \leq\exp\!\left(-\frac{\delta^2\mu}{2+\delta}\right).
\]
To prove simultaneous feasibility, the exponent must dominate the logarithm of the
number of constraints.  If independent rounding does not preserve a hard constraint,
use alteration, dependent rounding, or a problem-specific nonlinear map.

\subsection{Chapter-5 application pattern}

The later applications in the chapter repeatedly use the same accounting scheme:
\begin{enumerate}
\item write an LP whose optimum bounds $\OPT$ in the correct direction;
\item interpret $x$ as probabilities, possibly after scaling;
\item prove expected cost/value by linearity;
\item repair uncovered or overloaded constraints and charge the repair to the LP;
\item derandomize with conditional expectation or pessimistic estimators when needed.
\end{enumerate}
The technique, not merely the final constant, is the examinable invariant.

\section{SDP rounding and 3-colouring}

\subsection{Semidefinite programs (Section 6.1)}

A symmetric matrix $X$ is positive semidefinite, written $X\succeq0$, iff
$y^TXy\geq0$ for every $y$.  Equivalently, $X=V^TV$ and
$X_{ij}=v_i\cdot v_j$ for vectors $v_i$.  Thus an SDP can be viewed as a linear
optimization problem over a Gram matrix, or as a vector program.  The relaxation step
replaces products $x_ix_j$ of signs by inner products $v_i\cdot v_j$ of unit vectors.

\subsection{MAX-CUT and random hyperplanes (Section 6.2)}

With signs $x_i\in\{-1,1\}$, an edge $ij$ is cut by
$\tfrac12(1-x_ix_j)$.  Relax to unit vectors:
\[
 \max\ \frac12\sum_{ij\in E}w_{ij}(1-v_i\cdot v_j),
 \qquad \|v_i\|=1.
\]
Choose a random Gaussian vector $r$ and set $x_i=\operatorname{sign}(v_i\cdot r)$.
If the angle between $v_i,v_j$ is $\theta_{ij}$, then
\[
 \Prb[ij\text{ cut}]=\frac{\theta_{ij}}{\pi},
 \qquad v_i\cdot v_j=\cos\theta_{ij}.
\]
For every $\theta\in[0,\pi]$,
\[
 \frac{\theta}{\pi}\geq
 \alpha_{GW}\frac{1-\cos\theta}{2},
 \qquad
 \alpha_{GW}:=\min_{\theta\in[0,\pi]}
 \frac{2\theta}{\pi(1-\cos\theta)}\approx0.87856.
\]
Summing edgewise gives $\E[\ALG]\geq\alpha_{GW}\OPT_{\rm SDP}
\geq\alpha_{GW}\OPT$.

\subsection{Colouring a 3-colourable graph (Section 6.5)}

\textbf{Combinatorial baseline.}  While a vertex has degree at least $\sqrt n$, give
it one new colour and 2-colour its neighbourhood with two more new colours; the
neighbourhood is bipartite because the graph is 3-colourable.  Remove these vertices.
This occurs at most $\sqrt n$ times.  Greedily colour the remaining graph of maximum
degree below $\sqrt n$.  Total: at most $4\sqrt n$ colours.

For the SDP method solve
\[
 \min\lambda\quad\text{s.t. }v_i\cdot v_j\leq\lambda\ (ij\in E),
 \quad \|v_i\|=1.
\]
A proper 3-colouring gives three unit vectors at mutual angle $2\pi/3$, hence
$\lambda\leq-1/2$.  Choose
$t=2+\log_3\Delta$ independent random hyperplanes and colour a vertex by its
$t$-bit sign pattern.  Adjacent vertices share a pattern with probability at most
$(1/3)^t\leq1/(9\Delta)$.  Since $m\leq n\Delta/2$, Markov's inequality shows that
with probability at least $1/2$ at most $n/4$ edges are monochromatic.  Deleting one
endpoint from each such edge leaves at least $n/2$ correctly coloured vertices.
This is a semicolouring with
\[
  2^t=4\Delta^{\log_3 2}
\]
colours.  Repeating on the uncoloured half costs an extra factor $O(\log n)$.
Balancing the treatment of high-degree neighbourhoods with this semicolouring yields
\[
  \widetilde O\!\left(n^{\frac{\log_3 2}{1+\log_3 2}}\right)
  =\widetilde O(n^{0.387})
\]
colours.

\subsection{Large independent sets from the SDP (Section 13.2)}

Use the feasible vector system
\[
  v_i\cdot v_j=-\frac12\quad(ij\in E),
  \qquad \|v_i\|=1.
\]
Draw $r\sim N(0,I)$, let
\[
 \eta=\sqrt{\frac23\ln\Delta},\qquad
 S=\{i:v_i\cdot r\geq\eta\},
\]
and keep only vertices of $S$ with no neighbour in $S$, obtaining an independent set
$S'$.  If $ij\in E$ and both projections exceed $\eta$, orthogonal decomposition of
$v_j$ relative to $v_i$ implies an independent normal coordinate exceeds
$\sqrt3\eta$.  Hence
\[
 \Prb[j\in S\mid i\in S]\leq\overline\Phi(\sqrt3\eta),
 \qquad
 \Prb[i\notin S'\mid i\in S]\leq
 \Delta\overline\Phi(\sqrt3\eta)\leq\frac12.
\]
Using the Gaussian tail estimates
\[
 \frac{x}{1+x^2}\phi(x)\leq\overline\Phi(x)
 \leq\frac1x\phi(x),
\]
one obtains
\[
 \E|S'|=\Omega\!\left(n\Delta^{-1/3}(\ln\Delta)^{-1/2}\right).
\]
Repeated removal therefore colours with
$O(\Delta^{1/3}\sqrt{\ln\Delta}\log n)$ colours; combining regimes gives the
$\widetilde O(n^{1/4})$ bound discussed in the section/exercise.

\section{The primal--dual method}

\subsection{Set cover review (Section 7.1)}

For sets $S_j$ of weight $w_j$, the relaxation and dual are
\[
 \begin{array}{ll}
 \min & \sum_jw_jx_j\\
 \text{s.t.}&\sum_{j:e_i\in S_j}x_j\geq1\quad(i),\quad x_j\geq0,
 \end{array}
 \qquad
 \begin{array}{ll}
 \max & \sum_i y_i\\
 \text{s.t.}&\sum_{i:e_i\in S_j}y_i\leq w_j\quad(j),\quad y_i\geq0.
 \end{array}
\]
Raise a dual variable for an uncovered element until some set becomes tight, and take
that set.  If every element occurs in at most $f$ sets, then
\[
 \sum_{j\in I}w_j
 =\sum_{j\in I}\sum_{i\in S_j}y_i
 \leq f\sum_i y_i
 \leq f\OPT.
\]
The factor comes from the number of times each $y_i$ is charged, not from weak duality
alone.

\subsection{Shortest $s$--$t$ path (Section 7.3)}

The cut relaxation is
\[
 \min\sum_ec_ex_e
 \quad\text{s.t. }\sum_{e\in\delta(S)}x_e\geq1
 \quad(s\in S,t\notin S),\quad x_e\geq0.
\]
Its dual packs $s$--$t$ cuts:
\[
 \max\sum_Sy_S
 \quad\text{s.t. }\sum_{S:e\in\delta(S)}y_S\leq c_e,quad y_S\geq0.
\]
Grow the dual variable of the current component containing $s$ until an edge constraint
becomes tight, add the edge, and continue.  Tightness implies that the length assigned
to every reached vertex equals its shortest-path distance.  When $t$ is reached, the
tight predecessor edges recover an optimum path.  This is the exact primal--dual
prototype behind Dijkstra's algorithm.

\subsection{Generalized Steiner tree/Steiner forest (Section 7.4)}

Given pairs $(s_i,t_i)$, choose a minimum-cost forest connecting every pair.  A set
$S\subset V$ is active if it separates at least one still-unconnected demand pair.
The cut LP and dual are
\[
 \min\sum_ec_ex_e,quad
 \sum_{e\in\delta(S)}x_e\geq1\ (S\text{ separates a pair});
\]
\[
 \max\sum_Sy_S,quad
 \sum_{S:e\in\delta(S)}y_S\leq c_e.
\]
Start with singleton components and raise the dual variables of all active components
uniformly.  When an edge becomes tight, add it and merge its components.  Once all
pairs are connected, process added edges in reverse and delete every redundant edge.
For the pruned forest $F'$ and the active components $\mathcal C$ in any iteration,
\[
  \sum_{C\in\mathcal C}|\delta(C)\cap F'|\leq2|\mathcal C|.
\]
Integrating this inequality over all dual-growth intervals gives
\[
 c(F')\leq2\sum_Sy_S\leq2\OPT.
\]
The reverse-delete step is essential for the degree/charging inequality.

\section{Tree metric embeddings and buy-at-bulk}

\subsection{FRT embedding (Section 8.5)}

For a finite metric $(V,d)$, an FRT-type random hierarchical decomposition chooses a
random permutation of $V$ and a random scale shift, then partitions at geometrically
decreasing radii.  The cluster hierarchy defines a tree metric $d_T$.

The two required properties are
\[
 d_T(u,v)\geq d(u,v)\quad\text{for all }u,v,
 \qquad
 \E[d_T(u,v)]\leq O(\log n)d(u,v).
\]
Noncontraction is built into the tree-edge lengths: once $u$ and $v$ separate at scale
$i$, their tree path has length at least their original distance.  For expected stretch,
charge the first scale that separates them.  At scale $i$, a potential centre $w$ can
separate the pair only when the random radius falls in an interval of length at most
$d(u,v)$.  Among the centres close enough to matter, $w$ is the first in the random
permutation with probability equal to the reciprocal of its distance rank.  Summing
these probabilities produces $H_n=\sum_{j=1}^n1/j=O(\log n)$.  Multiplying the
separation probability at scale $i$ by the $O(2^i)$ tree distance charged at that scale,
and then summing over scales, yields
\[
  \E[d_T(u,v)]\leq O(H_n)d(u,v)=O(\log n)d(u,v).
\]

\textbf{Transfer rule.}  If a problem is solved exactly or within factor $\rho$ on a
tree and its cost is a nonnegative linear combination of pair distances, sampling the
tree and mapping back often gives expected factor $O(\rho\log n)$.

\subsection{Buy-at-bulk network design (Section 8.6)}

The input is a graph with lengths $\ell_e$, demand pairs $(s_i,t_i)$ with demands
$d_i$, and a nondecreasing concave cost function $f$ with $f(0)=0$.  Choose paths
$P_i$ and let $c_e=\sum_{i:e\in P_i}d_i$; minimize
\[
  \sum_{e}\ell_e f(c_e).
\]
On a tree the path for each pair is unique, so the optimum is immediate.  Sample a
dominating FRT tree, route on its unique paths, and replace every tree edge by a
shortest path in the original graph.  Domination makes the tree a valid upper model;
expected stretch contributes $O(\log n)$; concavity/subadditivity ensures that merging
mapped routes on an original edge does not increase the charged cost.  Therefore the
randomized algorithm is an $O(\log n)$-approximation.

\section{Planar maximum independent set PTAS (Section 10.2)}

Maximum independent set is extremely hard to approximate in general graphs, but planar
structure permits a PTAS.  Use a planar embedding and partition vertices into layers.
For an integer $k$, and each residue $r\in\{0,\ldots,k-1\}$, delete every layer whose
index is $r\pmod k$.  Each remaining component is $(k-1)$-outerplanar and has treewidth
$O(k)$, so weighted maximum independent set can be solved exactly by dynamic
programming for fixed $k$.

Let $O$ be an optimum independent set.  Across the $k$ choices of residue, every vertex
of $O$ is deleted exactly once.  Hence some residue deletes weight at most
$w(O)/k$.  Running all $k$ cases and taking the best gives
\[
  \ALG\geq\left(1-\frac1k\right)\OPT.
\]
Choose $k\geq1/\varepsilon$ for a $(1-\varepsilon)$-approximation.  Planarity is used
to turn deletion of one layer class into bounded treewidth; the averaging argument
alone is not enough.

\section{Hardness of approximation}

\subsection{Direct gap reductions (Section 16.1)}

A gap reduction maps an NP-complete instance $x$ to an optimization instance $I_x$ so
that
\[
 x\in L\Rightarrow\OPT(I_x)\leq a,\qquad
 x\notin L\Rightarrow\OPT(I_x)\geq b
\]
for minimization (reverse the inequalities for maximization).  Any ratio strictly
better than $b/a$ distinguishes the cases.

The section's main example reduces 3-dimensional matching to makespan scheduling on
unrelated parallel machines.  The construction creates element jobs, triple machines,
and dummy jobs so that a perfect 3-dimensional matching yields a schedule of length
two, whereas every schedule for a NO instance has length at least three.  Consequently,
unless $P=NP$, there is no $(3/2-\varepsilon)$-approximation for any fixed
$\varepsilon>0$.  The proof must establish both directions of the gap; mere
NP-hardness of computing $\OPT$ gives no approximation lower bound.

\subsection{PCP reductions (Section 16.3)}

$L\in\mathrm{PCP}_{c,s}(r(n),q(n))$ means that a randomized polynomial-time verifier
uses at most $r(n)$ random bits, reads at most $q(n)$ locations of a proof $\pi$, and
satisfies
\[
 x\in L\Rightarrow\exists\pi:\Pr[V^\pi(x)=1]\geq c,qquad
 x\notin L\Rightarrow\forall\pi:\Pr[V^\pi(x)=1]\leq s.
\]
The PCP theorem gives
\[
  \mathsf{NP}=\mathrm{PCP}(O(\log n),O(1))
\]
with a constant completeness--soundness gap.

To obtain an optimization instance, create one constant-arity constraint for every
verifier random string; its variables are the proof locations queried on that string.
If the verifier has perfect completeness, a YES instance satisfies all constraints.
For a NO instance, every assignment satisfies at most an $s$ fraction.  Thus a PCP is
exactly a source of a robust gap rather than merely an NP witness verifier.

The refined three-query construction in the section yields MAX-E3SAT hardness:
for every constant $\delta>0$ it is NP-hard to distinguish a satisfiable instance from
one in which at most a $(7/8+\delta)$ fraction of clauses can be satisfied.  Hence, for
any constant $\alpha>7/8$, an $\alpha$-approximation for MAX-E3SAT would imply
$P=NP$.  This matches the $7/8$ random-assignment algorithm up to an arbitrarily small
additive constant.

\section*{Primary reference}
D.~P.~Williamson and D.~B.~Shmoys,
\emph{The Design of Approximation Algorithms}, Cambridge University Press, 2011,
electronic web edition.

\end{document}
